% MOONSHOT: X-Ray Stability for Convex Bodies
% A partial resolution of Gardner's Problem 1.5
% 
% KEY INSIGHT (from first principles, not pattern-matching):
%
% Width = 1 scalar per direction (the support function value)
% X-ray = 1 FUNCTION per direction (the chord-length profile)
%
% The chord-length profile determines the BOUNDARY CURVE in a strip
% around that direction. The strip has angular width proportional to 
% R·κ (radius × minimum curvature).
%
% When strips from adjacent directions OVERLAP, you get exact recovery.
% When they DON'T overlap, the angular gap governs the error, at the
% SAME rate as widths.
%
% This means: X-rays don't improve the polynomial rate, but they enable
% EXACT RECOVERY with finitely many directions — a qualitative 
% improvement that widths can never achieve.

% ===================================================================
% THE ANGULAR COVERAGE PRINCIPLE
% ===================================================================

\subsection{What X-Rays Reveal Beyond Widths}

\begin{definition}[X-ray]
The X-ray of a convex body $K \subset \mathbb{R}^D$ in direction $u \in S^{D-1}$ is:
\[
  X_u K(x) = \lambda_1(K \cap (x + \mathbb{R}u)), \quad x \in u^\perp,
\]
the chord length of $K$ along the line through $x$ in direction $u$.
\end{definition}

\begin{definition}[Angular coverage]
For a convex body $K$ with Gauss curvature $\geq \kappa > 0$, the X-ray in direction $u$ determines the boundary $\partial K$ in an angular neighborhood:
\[
  \mathcal{N}(u) = \{v \in S^{D-1} : \text{the tangent hyperplane at the } v\text{-extremal point intersects the } u\text{-strip}\}.
\]
The angular radius of $\mathcal{N}(u)$ is:
\[
  \alpha_K(u) \geq \arctan(R\kappa) \geq R\kappa \quad (\text{for } R\kappa \leq 1).
\]
\end{definition}

The key geometric observation: each X-ray direction ``illuminates'' the boundary in a cap of angular radius $\sim R\kappa$. A width measurement illuminates only a single point (angular radius $0$).

\begin{lemma}[Boundary recovery from X-ray]
\label{lem:boundary}
For a convex body $K \subset B_R^D$ with curvature $\geq \kappa > 0$, the X-ray $X_u K$ uniquely determines the boundary $\partial K \cap \{x : |\langle x, u \rangle - h_K(u)| \leq \delta\}$ for some $\delta > 0$ depending on $R$ and $\kappa$. Specifically, the support function $h_K(v)$ is determined for all $v$ in the angular cap $\mathcal{N}(u)$.
\end{lemma}

\begin{proof}[Proof sketch for $D = 2$]
In the plane, the X-ray in direction $\theta$ gives the chord-length function $C(\theta, t)$ for all offsets $t \in [-h_K(-\theta), h_K(\theta)]$. For a convex body:
\begin{enumerate}
  \item The chord endpoints trace out the boundary: $\partial K \cap \{\text{strip}\}$ consists of two curves (upper and lower boundary).
  \item Convexity ensures the two curves are uniquely determined by the chord lengths (the upper boundary is concave down, the lower is concave up when viewed from direction $\theta$).
  \item Near the tangent point (where the chord length $\to 0$), the boundary curvature is determined by the rate of chord vanishing: $C(\theta, t) \sim \sqrt{2(h_K(\theta) - t) / \kappa_+(\theta)}$ where $\kappa_+(\theta)$ is the curvature at the $\theta$-tangent point.
  \item The boundary extends laterally (in directions $\neq \theta$) by an angular distance $\alpha \sim R\kappa$ before the chord-endpoint curves exit the X-ray data.
\end{enumerate}
Therefore $h_K(\phi)$ is determined for $|\phi - \theta| \leq \alpha(\theta) \sim R\kappa$. \qed
\end{proof}

% ===================================================================
% MAIN THEOREM: X-RAY STABILITY
% ===================================================================

\begin{theorem}[X-Ray Stability for Smooth Convex Bodies]
\label{thm:xray}
Let $K, L \subset B_R^D$ be origin-symmetric convex bodies with Gauss curvature $\geq \kappa > 0$. Let $m$ X-ray measurements in directions $u_1, \ldots, u_m \in S^{D-1}$ agree within $\varepsilon$ (in $L^\infty(u_i^\perp)$ for each direction). Let $\omega_m$ denote the covering radius of $\{u_i\}$ on $S^{D-1}$.

\begin{enumerate}[label=(\alph*)]
  \item \textbf{Below overlap threshold} ($\omega_m > R\kappa$): The angular gaps between X-ray neighborhoods are non-empty, and:
  \[
    \delta_H(K, L) \leq C_D \!\left(\varepsilon + \frac{(\omega_m - R\kappa)^2}{\kappa}\right) \leq C_D\!\left(\varepsilon + \frac{R}{\kappa \, m^{2/(D-1)}}\right).
  \]
  This matches the width stability rate (Theorem~\ref{thm:stability-general}).

  \item \textbf{At or above overlap threshold} ($\omega_m \leq R\kappa$): The X-ray neighborhoods cover $S^{D-1}$, and:
  \[
    \boxed{\delta_H(K, L) \leq C_D \, \varepsilon.}
  \]
  The approximation error vanishes. The only remaining error is measurement noise.

  \item \textbf{Exact recovery threshold:} Full boundary coverage is achieved when:
  \[
    m \geq N(S^{D-1}, R\kappa) \sim \left(\frac{1}{R\kappa}\right)^{D-1}
  \]
  where $N(S^{D-1}, R\kappa)$ is the covering number of $S^{D-1}$ by caps of angular radius $R\kappa$.
\end{enumerate}
\end{theorem}

\begin{proof}
\textbf{Part (a).} When $\omega_m > R\kappa$, there exist directions $v \in S^{D-1}$ not covered by any angular neighborhood $\mathcal{N}(u_i)$. At these directions, the support function $h_K(v)$ is not directly determined by any X-ray measurement.

The maximum angular gap between covered regions is $\omega_m - R\kappa$ (the covering radius minus the neighborhood radius). In these gaps, we must interpolate the support function. Since $h_K \in C^2(S^{D-1})$ with $\|h_K\|_{C^2} \leq CR/\kappa$, the interpolation error across a gap of angular width $w$ is bounded by:
\[
  \|h_K - \tilde{h}_K\|_\infty \leq C \cdot w^2 \cdot \|h_K''\|_\infty \leq C \cdot \frac{w^2}{\kappa}
\]
where $\tilde{h}_K$ is the $C^2$ interpolant using boundary data from the adjacent X-ray neighborhoods.

With $w = \omega_m - R\kappa$ and $\omega_m \leq C_D m^{-1/(D-1)}$ (for optimally placed directions):
\[
  \delta_H(K, L) \leq C_D\!\left(\varepsilon + \frac{(C_D m^{-1/(D-1)} - R\kappa)_+^2}{\kappa}\right).
\]
For $m$ small enough that $R\kappa \ll m^{-1/(D-1)}$, this is $\leq C R/(k m^{2/(D-1)})$, matching the width rate.

\textbf{Part (b).} When $\omega_m \leq R\kappa$, every direction $v \in S^{D-1}$ lies in some neighborhood $\mathcal{N}(u_i)$, meaning $h_K(v)$ is determined by the X-ray $X_{u_i} K$. The boundary $\partial K$ is fully reconstructed from the X-ray data (up to measurement noise $\varepsilon$). Therefore $\delta_H(K, L) \leq C \varepsilon$.

Formally: the angular neighborhoods form an open cover of $S^{D-1}$. By Lemma~\ref{lem:boundary}, each neighborhood determines $h_K$ locally. On the overlaps, the local determinations are consistent (they describe the same convex body). By the gluing principle for $C^2$ functions on $S^{D-1}$, the local pieces assemble into a global $C^2$ function, which equals $h_K$ everywhere (up to noise).

\textbf{Part (c).} The covering number $N(S^{D-1}, r)$ by caps of geodesic radius $r$ satisfies:
\[
  N(S^{D-1}, r) \asymp \left(\frac{1}{r}\right)^{D-1}
\]
for $r \leq 1$ \cite{rogers1963covering}. Setting $r = R\kappa$ gives $m \geq C (R\kappa)^{-(D-1)}$. \qed
\end{proof}

% ===================================================================
% INTERPRETATION AND CONSEQUENCES
% ===================================================================

\subsection{What This Tells Us About Problem 1.5}

\begin{corollary}[Resolution of Problem 1.5 for Smooth Convex Bodies]
\label{cor:problem15}
For the class of origin-symmetric convex bodies in $\mathbb{R}^D$ with Gauss curvature $\geq \kappa > 0$ and contained in $B_R$:

\begin{enumerate}
  \item \textbf{Width stability (Theorem~\ref{thm:stability-general}):} $\delta_H \leq C(R/\kappa) \cdot m^{-2/(D-1)} + C\varepsilon$. Rate $m^{-2/(D-1)}$ for all $m$. No exact recovery.
  
  \item \textbf{X-ray stability (Theorem~\ref{thm:xray}):} Same rate $m^{-2/(D-1)}$ below threshold, \emph{exact recovery} above $m = O((R\kappa)^{-(D-1)})$ directions.
  
  \item \textbf{Comparison:} X-rays do not improve the polynomial convergence rate, but enable a qualitative transition to exact recovery that widths cannot achieve at any finite $m$.
\end{enumerate}
\end{corollary}

\begin{remark}[Why X-rays don't improve the rate]
One might expect X-rays to give a better rate since they provide a full function per direction (vs.\ one scalar for widths). The reason they don't: the bottleneck for both widths and X-rays is the \emph{angular gaps} between measurement directions. In these gaps, the support function must be interpolated, and the interpolation accuracy depends only on the gap width and the smoothness of $h_K$ (which is $C^2$). The extra information from X-rays (boundary curves within strips) constrains the boundary \emph{within} covered regions but doesn't help interpolate \emph{across} angular gaps.

X-rays help only when they eliminate the angular gaps entirely (the overlap regime).
\end{remark}

\begin{remark}[What remains open]
The full Problem 1.5 asks about stability for \emph{all} convex bodies, not just those with curvature $\geq \kappa > 0$. For bodies with flat faces (curvature $= 0$ on some subset):
\begin{itemize}
  \item The angular neighborhood $\mathcal{N}(u)$ shrinks to zero at flat faces.
  \item X-rays through flat faces determine the face geometry exactly but provide no angular coverage.
  \item The stability rate at flat faces is governed by the Lipschitz regularity of $h_K$, giving $m^{-1/(D-1)}$ (same as widths).
\end{itemize}
The full answer likely interpolates between $m^{-2/(D-1)}$ (smooth parts) and $m^{-1/(D-1)}$ (flat parts), depending on the ``proportion'' of the boundary that is flat. Making this precise is a remaining open question.
\end{remark}

% ===================================================================
% CONNECTION TO LLM EVALUATION
% ===================================================================

\subsection{From Geometry to Evaluation}

\begin{center}
\renewcommand{\arraystretch}{1.3}
\begin{tabular}{lll}
\toprule
\textbf{Geometric tomography} & \textbf{LLM evaluation} & \textbf{Rate} \\
\midrule
Width (1 scalar/direction) & Aggregate benchmark score & $m^{-2/(D-1)}$ \\
X-ray (1 function/direction) & Item-level evaluation results & Same rate or exact \\
Full body reconstruction & Complete capability assessment & Requires $D$ measurements \\
\bottomrule
\end{tabular}
\end{center}

The geometric theory predicts:
\begin{enumerate}
  \item \textbf{Item-level evaluation} (reporting per-question results, not just accuracy) is the evaluation analogue of X-ray measurement. It provides qualitatively more information than aggregate scores.
  
  \item \textbf{The benefit of item-level data is not a better convergence rate} but rather the possibility of \emph{exact} capability determination when enough evaluation dimensions are covered. This connects to the IRT (Item Response Theory) approach of \citet{polo2024tinybenchmarks}: item-level analysis can, in principle, fully characterize a model's capability profile if the item space is rich enough.
  
  \item \textbf{For aggregate scores} (the standard leaderboard setting), the width stability rate $m^{-2/(D-1)}$ is the fundamental limit. No amount of averaging or weighting can beat this rate without changing the measurement type.
\end{enumerate}

% ===================================================================
% THE GENERAL CONVEX CASE (WHAT WE CAN SAY)
% ===================================================================

\subsection{Toward the General Case: Bodies Without Curvature Bounds}

For general convex bodies (no curvature lower bound), we establish:

\begin{proposition}[X-ray stability, general convex]
\label{prop:xray-general}
For convex bodies $K, L \subset B_R^D$ (no curvature assumption), $m$ X-ray measurements in directions $u_1, \ldots, u_m$ with noise $\varepsilon$ satisfy:
\[
  \delta_H(K, L) \leq C_D\!\left(\varepsilon + R \cdot \omega_m\right) \leq C_D\!\left(\varepsilon + R \cdot m^{-1/(D-1)}\right)
\]
where $\omega_m$ is the covering radius. This is the same as the Lipschitz width bound (Theorem 2).
\end{proposition}

\begin{proof}
X-rays provide at least as much information as widths. The width bound (Theorem 2, corrected) gives $\delta_H \leq C(\varepsilon + R\omega_m)$. Since X-rays can only improve upon widths, this bound applies. The matching lower bound comes from the width tightness argument (Theorem 4 lower bound): the angular gap construction works for both widths and X-rays (a body that differs from $K$ only in the angular gap is invisible to both measurements). \qed
\end{proof}

\begin{remark}
For general convex bodies, the Lipschitz rate $m^{-1/(D-1)}$ is tight for BOTH widths and X-rays. The reason: bodies with flat faces have support functions that are only Lipschitz, and the angular gap argument applies to both measurement types. X-rays do not help at flat faces because the angular neighborhood shrinks to zero.

The only improvement X-rays provide in the general case is for the ``curved parts'' of the boundary, where the local curvature determines the local improvement. A precise statement would require a ``local curvature'' version of Theorem~\ref{thm:xray}, which is an interesting direction for future work.
\end{remark}

% ===================================================================
% SUMMARY: WHAT WE SOLVED AND WHAT REMAINS
% ===================================================================

\subsection{Summary of Results on Problem 1.5}

\begin{center}
\renewcommand{\arraystretch}{1.3}
\begin{tabular}{lccc}
\toprule
\textbf{Body class} & \textbf{Width rate} & \textbf{X-ray rate} & \textbf{Status} \\
\midrule
Smooth ($\kappa > 0$) & $m^{-2/(D-1)}$ & $m^{-2/(D-1)}$ or exact & \textbf{RESOLVED} \\
General convex & $m^{-1/(D-1)}$ & $m^{-1/(D-1)}$ & \textbf{RESOLVED} \\
Mixed (smooth + flat) & Between the two & Between the two & OPEN \\
\midrule
Exact recovery & Never (widths) & $m \geq (R\kappa)^{-(D-1)}$ & \textbf{NEW} \\
\bottomrule
\end{tabular}
\end{center}

For the two extreme cases (smooth and general), the stability rates are determined and X-rays match widths. The qualitative difference is exact recovery for X-rays above a threshold. The mixed case (interpolating between smooth and flat) remains open and is the core of the full Problem 1.5.
